
% -------------------- 符号说明表 -------------------- %
\begin{center}
\heiti\zihao{-3}符号说明表
\end{center}%

该部分内容是对于论文中所使用的主要符号、标志、缩略语、首字母缩写、计量单位、自定义名词和术语等的注释说明，应进行统一整理，用于清晰的展现出论文的符号等的具体含义，有助于论文的审阅。若上述符号使用数量不多，可以不设此部分，但必须在论文中出现时加以说明。

下面以物理学中的一些符号为例进行示范：

% 跨页表格
{\setlength{\tabcolsep}{15mm}%设置列间距
\begin{longtable}{ccc} %可在此使用p{数值}设置列宽（注：p{...}默认居左表示）

    %\caption{The List of Notation} \label{tab:Notiations} \\
	\toprule
	{\bf 符~~号} & {\bf 代表意义} & {\bf 单~~位}\\
	%\midrule
    \endfirsthead %以上是首页表头
    \midrule
    {\bf 符号} & {\bf 解释} & {\bf 单位}\\
    \midrule
    \endhead %以上是通用表头
    \midrule
    \endfoot %以上是通用表尾
    \bottomrule
    \endlastfoot %以上是末页表尾，以下是表格正文
    $A$ & 截面积；散热面积 & \unit{m^2} \\
    $F$ & 力 & \unit{N} \\
    $e$ & 电子电荷 & \unit{V} \\
    $I, i$ & 电流 & \unit{A} \\
    $m$ & 质量 & \unit{Kg} \\
    $Q$ & 热量或流量 & \unit{Ki} \\
    $H$ & 焓值 & \unit{Ki/Kg} \\
    $S$ & 熵值 & \unit{Ki/Kg} \\
\end{longtable}}%
